Eine feste Menge eines idealen Gases ist im Zustand A bei $p_A = p_0$, $V_A = 2\,V_0$ und $T_A = 2\,T_0$. Es durchläuft folgenden Kreisprozess:
\begin{itemize}
\item A $\to$ B: isochor, bis sich der Druck verdreifacht hat;
\item B $\to$ C: isotherme Kompression auf das halbe Volumen;
\item C $\to$ D: isochor, bis der Druck auf einen Drittel gesunken ist;
\item D $\to$ A: isotherm zurück in den Anfangszustand.
\end{itemize}
Die Achsen sind in Einheiten von $p_0$, $V_0$ und $T_0 = \frac{p_0V_0}{nR}$ beschriftet; damit ist $T/T_0 = (p/p_0)\cdot(V/V_0)$. Ein Schüler skizziert den Prozess im $p(T)$-Diagramm:

\begin{center}
\begin{tikzpicture}[x=0.75cm, y=0.75cm, >=Stealth, font=\small]
 \draw[black!12] (0,0) grid[xstep=1, ystep=1] (7.001,7.001);
 \draw[->] (0,0) -- (7.6,0) node[right] {$T$ in $T_0$};
 \draw[->] (0,0) -- (0,7.6) node[above] {$p$ in $p_0$};
 \node[below left] at (0,0) {\small 0};
 \draw (1,0.08) -- (1,-0.08) node[below] {\small 1};
 \draw (2,0.08) -- (2,-0.08) node[below] {\small 2};
 \draw (3,0.08) -- (3,-0.08) node[below] {\small 3};
 \draw (4,0.08) -- (4,-0.08) node[below] {\small 4};
 \draw (5,0.08) -- (5,-0.08) node[below] {\small 5};
 \draw (6,0.08) -- (6,-0.08) node[below] {\small 6};
 \draw (7,0.08) -- (7,-0.08) node[below] {\small 7};
 \draw (0.08,1) -- (-0.08,1) node[left] {\small 1};
 \draw (0.08,2) -- (-0.08,2) node[left] {\small 2};
 \draw (0.08,3) -- (-0.08,3) node[left] {\small 3};
 \draw (0.08,4) -- (-0.08,4) node[left] {\small 4};
 \draw (0.08,5) -- (-0.08,5) node[left] {\small 5};
 \draw (0.08,6) -- (-0.08,6) node[left] {\small 6};
 \draw (0.08,7) -- (-0.08,7) node[left] {\small 7};
 \draw[black!12] (0,0) grid[xstep=1, ystep=0.5] (7.001,7.001);
 \draw[very thick, blue] (2,1) -- (6,3);
 \draw[-{Stealth[length=2.6mm]}, blue] (4.191,2.096) -- (4.209,2.104);
 \draw[very thick, blue] (6,3) -- (6,1.5);
 \draw[-{Stealth[length=2.6mm]}, blue] (6,2.185) -- (6,2.165);
 \draw[very thick, blue] (6,1.5) -- (2,0.5);
 \draw[-{Stealth[length=2.6mm]}, blue] (3.81,0.952) -- (3.79,0.948);
 \draw[very thick, blue] (2,0.5) -- (2,1);
 \draw[-{Stealth[length=2.6mm]}, blue] (2,0.765) -- (2,0.785);
 \fill[black] (2,1) circle (2pt) node[above left] {$A$};
 \fill[black] (6,3) circle (2pt) node[above right] {$B$};
 \fill[black] (6,1.5) circle (2pt) node[right] {$C$};
 \fill[black] (2,0.5) circle (2pt) node[below left] {$D$};
\end{tikzpicture}
\end{center}

\begin{abcliste}
\abc Bestimmen Sie $p$, $V$ und $T$ in den Zuständen A bis D. Schliesst sich der Kreisprozess?
\abc In der Skizze ist ein Schritt falsch. Welcher? Ist die Art des Prozesses falsch, seine Form im Diagramm oder der Faktor?
\abc Die Schritte C $\to$ D und D $\to$ A haben in der Skizze die richtige Form. Warum liegen C und D trotzdem falsch?
\abc Skizzieren Sie das richtige $p(T)$-Diagramm.
\end{abcliste}
